\section{Tolerant B-reps and realizations}\label{sec:setting}

\subsection{Complexes}
$K$ is a finite cell complex in the B-rep sense. Its 0-cells are points, its edges are arcs or closed curves, and
its faces are compact surfaces, possibly with holes, whose interiors are the open faces. The underlying space $|K|$ is a
disjoint union of closed connected surfaces $|K_1|,\dots,|K_n|$ (the \emph{shells}). We write $\tau\le\sigma$ when $\tau$ is a face of $\sigma$ (including
$\tau=\sigma$). For the uniqueness theorem we assume $K$ is \emph{regular}: every closed cell $\bar\sigma$ is a
ball whose boundary is a subcomplex. B-rep complexes with seam edges, closed edges or faces with holes are not
regular, and must first be subdivided. A new cell inherits the zone of the open cell of $K$ that contains it.
This keeps the zones monotone, and a realization of the original model is a realization of the subdivision, since
$h(\inter\sigma')\subseteq h(\inter\sigma)\subseteq\Z(\sigma)$ for a new cell $\sigma'\subseteq\sigma$. The carrier
condition of Section~\ref{sec:uniqueness} must then hold for the subdivided complex. A new edge inside a face $f$
has carrier $\Z(f)$, so subdividing faces adds the hypothesis $\pi_1(\Z(f))=0$; contractible zones satisfy it. A map of $|K|$ is \emph{PL} if it is linear on the simplices of some triangulation of $|K|$ that subdivides its
cells. In dimension two, any two such triangulations are related by a cell-preserving PL homeomorphism (standard; not
separately checked against a source), so the notion does not depend on the choice. New vertices created inside an edge or face inherit its zone, so the carrier condition requires those
zones to be path-connected. The existence results say where they apply to B-rep
complexes directly.

\subsection{Geometric data and zones}
\begin{itemize}
\item Each vertex $v$ has a point $p_v$ and a tolerance $\tau_v$.
\item Each edge $e$ has a $C^2$ curve $\gamma_e\colon[0,1]\to\R^3$ and a tolerance $\tau_e$.
\item Each face $f$ has a $C^2$ surface $S_f$, a compact trimmed patch $P_f\subseteq S_f$, and a tolerance
$\tau_f$. The boundary of $P_f$ consists of the images $\beta_{f,e}=S_f\circ(\text{pcurve})$ of the pcurves of the
edges $e<f$.
\end{itemize}
The \emph{deviation} of $e$ is $\mathrm{dev}_e=\max_{f>e}\max_t|\beta_{f,e}(t)-\gamma_e(t)|$, measured at equal
parameters. OpenCascade's SameParameter property asks for $\mathrm{dev}_e\le\tau_e$ in this sense; its format
specification defines the edge tolerance as a one-sided distance from the pcurve images to the 3D
curve~\cite{OCCT}.

\begin{definition}[Zones]\label{def:zones}
We use \emph{augmented} open zones:
\begin{itemize}
\item $\Z(v)=B(p_v,\tau_v)$;
\item $\Z(e)=N_{\tau_e}(\gamma_e)\cup\bigcup_{v<e}\Z(v)$;
\item $\Z(f)=N_{\tau_f}(P_f)\cup\bigcup_{e<f}\Z(e)$.
\end{itemize}
Here $N_r$ is the open $r$-neighbourhood. Augmentation makes $\Z$ monotone: $\tau\le\sigma$ implies
$\Z(\tau)\subseteq\Z(\sigma)$.
\end{definition}

\begin{definition}[Realization]\label{def:realization}
A \emph{realization} of the tolerant B-rep $T=(K,\text{data},\tau)$ is a PL embedding $h\colon|K|\to\R^3$ with
$h(\inter\sigma)\subseteq\Z(\sigma)$ for every cell $\sigma$.
\end{definition}

\begin{remark}\label{rem:opencell}
\begin{itemize}
\item \emph{Open-cell condition.} We require only $h(\inter\sigma)\subseteq\Z(\sigma)$. The closed-cell
condition with unaugmented neighbourhoods, $h(\bar\sigma)\subseteq N_{\tau_\sigma}(\cdot)$, usually fails for every
candidate, mainly because face tolerances are tiny: $h(v)$ would have to lie within $\tau_f$ of every incident patch. With augmented zones the open-cell
condition implies $h(\bar\sigma)\subseteq\Z(\sigma)$, since $\bar\sigma$ is the union of the open cells
$\tau\le\sigma$.
\item \emph{Tameness.} Requiring $h$ to be PL rules out wild embeddings such as Alexander's horned sphere.
\end{itemize}
\end{remark}

\begin{definition}[Carriers]
The \emph{carrier} of a cell is $\C(\sigma)=\bigcup_{\tau\le\sigma}\Z(\tau)$. For augmented zones
$\C(\sigma)=\Z(\sigma)$. A continuous map $g\colon|K|\to\R^3$ is \emph{carried} if $g(\bar\sigma)\subseteq
\C(\sigma)$ for every $\sigma$. By Remark~\ref{rem:opencell}, every realization is carried.
\end{definition}

\begin{definition}[Solid]\label{def:solid}
Let $S_1,\dots,S_n\subset\R^3$ be pairwise disjoint closed surfaces. The \emph{depth} of a point $x\notin\bigcup S_j$
is the number of $j$ for which $x$ lies in the bounded component of $\R^3\setminus S_j$. The solid bounded by
the $S_j$ is the closure of the set of points of odd depth.
\end{definition}
This matches the usual B-rep convention of outer shells and cavities when the shells are consistently oriented;
a consistency check of the orientation is separate.
